\begin{annexitem}{annex:p03.general-fibres}{thm:p03-general-fibres}{Uniform model fibres}
\begin{proof}
Choose fresh symbols as follows: a binary function \(p\); an
\((r+1)\)-ary relation \(R^\sharp\) for each \(r\)-ary relation
\(R\in\tau\); an \((r+1)\)-ary function \(F_f\) for each
\(r\)-ary function \(f\in\tau\); and a unary function \(C_c\) for
each constant \(c\in\tau\). Let \(\sigma\) contain these auxiliary symbols and the same single
sort as \(\tau\), and let \(d\) be a fresh constant. All symbols are chosen distinctly.
No cardinality bound on \(\tau\) is used in making this set of symbols.

For a \(\sigma\)-structure \(\mathfrak B\) and \(b\in B\), define a
\(\tau\)-structure \(\mathfrak A_b\) on the whole set \(B\) by
\[
\begin{aligned}
 R^{\mathfrak A_b}(\bar a)
   &\quad\Longleftrightarrow\quad
       (R^\sharp)^{\mathfrak B}(b,\bar a),\\
 f^{\mathfrak A_b}(\bar a)
   &=p^{\mathfrak B}\bigl(b,F_f^{\mathfrak B}(b,\bar a)\bigr),\\
 c^{\mathfrak A_b}
   &=p^{\mathfrak B}\bigl(b,C_c^{\mathfrak B}(b)\bigr).
\end{aligned}
\]
For a nullary relation the first line means simply that its truth value
is the truth value of \(R^\sharp(b)\). Every displayed function is total.
The image \(D_b\) of \(p^{\mathfrak B}(b,-)\) is nonempty because
\(B\) is nonempty; it contains every displayed constant and every
value of every displayed function. Thus it is \(\tau\)-closed in
\(\mathfrak A_b\). Put
\(\mathfrak B_b=\mathfrak A_b\,|\,D_b\).
Neither idempotence of \(p(b,-)\) nor any property of its restriction
to its image is needed for this conclusion.

We now obtain the claimed sentence using precisely the two semantic
operations of regularity. First relativize \(\varphi\) to the
first-order definable domain
\[
 \chi(y;d)\ :=\ \exists z\,\bigl(p(d,z)=y\bigr).
\]
Here \(y\) is the distinguished fresh constant used in the
relativization definition, and \(d\) remains a parameter constant.
The result is a sentence \(\varphi^{D_d}\) over
\(\tau\cup\{p,d\}\). Its defining satisfaction equivalence holds
whenever that domain is \(\tau\)-closed. Next, in that sentence,
simultaneously substitute all relation, function and constant symbols
of \(\tau\), retaining the sort: replace a
relation by \(R^\sharp(d,\bar y)\), a function by the graph formula
\[
 u=p\bigl(d,F_f(d,\bar y)\bigr),
\]
and a constant by the graph formula \(u=p(d,C_c(d))\).
The remaining symbols \(p,d\) are unchanged. Formally one first regards
\(\varphi^{D_d}\) as a sentence over
\(\tau\cup\sigma\cup\{d\}\), using vocabulary extension, and applies
simultaneous substitution with base vocabulary \(\sigma\cup\{d\}\).
The reduct property ensures that unused auxiliary symbols have no
other effect on its satisfaction.

Call the resulting sentence \(\widehat\varphi(d)\). The graph formulas
always define total functions and unique constants, so the expansion
required by the substitution property always exists and is exactly
\(\mathfrak A_b\), together with the unchanged auxiliary symbols.
The closed-domain condition for relativization was proved above.
Consequently substitution followed by the relativization equivalence
gives
\[
 (\mathfrak B,b)\models\widehat\varphi(d)
 \quad\Longleftrightarrow\quad
 \mathfrak A_b\,|\,D_b\models\varphi.
\]
This argument uses the satisfaction clauses for the two operations,
not an induction on a presumed syntax of \(\mathscr L\). In particular,
it applies when a sentence uses infinitely many vocabulary symbols.
The substitution is simultaneous even in that case; no collection of
additional graph axioms is added to the theory.

For the final assertion, let \((\mathfrak M_j)_{j\in J}\) be
any set-indexed family. The empty family imposes no requirement;
choose a nonempty ambient domain containing any prescribed extra
points and interpret the symbols arbitrarily. Otherwise replace
the domains by pairwise disjoint isomorphic copies \(M_j\),
also disjoint from distinct new index elements \(i_j\) and from
a prescribed set \(X\) of extra points. Choose \(a_j\in M_j\)
and take
\[
 B=X\ \mathbin{\dot\cup}\ \{i_j:j\in J\}
       \ \mathbin{\dot\cup}\ \coprod_{j\in J}M_j.
\]
Set \(p(i_j,z)=z\) for \(z\in M_j\), and
\(p(i_j,z)=a_j\) otherwise. On first arguments not among the
\(i_j\), set \(p\) to any fixed element of \(B\).
For arguments in \(M_j\), interpret \(F_f(i_j,\bar a)\)
as \(f^{\mathfrak M_j}(\bar a)\), interpret \(C_c(i_j)\)
as \(c^{\mathfrak M_j}\), and make
\(R^\sharp(i_j,\bar a)\) agree with
\(R^{\mathfrak M_j}(\bar a)\). This includes the empty tuple
for nullary relations. Complete other function values by a fixed
element of \(B\), and other relation values by false.
The prescriptions are compatible because the indices are distinct.
Then \(D_{i_j}=M_j\), the projection is the identity there,
and its induced structure is \(\mathfrak M_j\). The additional
points in \(X\) do not change any fibre, since their projections
at \(i_j\) are the chosen defaults \(a_j\).
\end{proof}
\end{annexitem}

\begin{annexitem}{annex:p03.general-pinning}{thm:p03-general-pinning}{Countable compactness and well-ordering numbers}
\begin{proof}
Regularity supplies equivalents of every first-order sentence, and
vocabulary extension and the reduct property preserve their meanings.
We use such sentences without changing their first-order notation.
A displayed free variable in an \(\mathscr L\)-expression denotes a
fresh parameter constant; its existential quantification is
particularization, and universal quantification is negated
particularization of a negation. Thus the quantified expressions below
are legitimate applications of the stated closure properties.

\emph{Countable compactness forbids infinite pinned ordinals.}
Suppose a countable theory \(T\) pins an infinite ordinal via \(<\),
and let \(\mathfrak A\models T\) witness an infinite order type.
Add fresh constants \(d_n\), for \(n<\omega\), and the sentences
\(d_{n+1}<d_n\). Every finite subset of this enlarged theory has a
model: keep \(\mathfrak A\), choose a finite descending chain long
enough for all constant indices occurring in the subset, assign these
constants to the chain, and assign the remaining constants arbitrarily.
An infinite well-order contains chains of every finite length.
Countable compactness would give a model of the entire enlarged theory.
Its reduct would satisfy \(T\), but \(\{d_n:n<\omega\}\) would be
a nonempty subset of its ordered field with no least element. This
contradicts pinning. Hence (1) implies (2).

For every finite \(n\geq2\), a first-order sentence says that \(<\)
is a strict linear order of its field and that its field has exactly
\(n\) elements. Every model has order type \(n\), and the usual
\(n\)-element order is a model. These single sentences show
\(w_{\aleph_0}(\mathscr L)\geq\omega\). Under (2), all the pinned
ordinals are finite, so the reverse inequality holds. Thus (2) implies (3).
This uses the field convention exactly: no singleton field of a strict
binary relation is presumed to exist.

\emph{Failure of compactness constructs pinned infinite orders.}
Suppose \(T\subseteq\mathscr L[\tau]\) is countable, each finite subset
has a model, and \(T\) itself has none. Such a theory cannot be finite.
Write it as \(\{\varphi_n:n<\omega\}\). Apply
\cref{thm:p03-general-fibres} once to \(\tau\), obtaining a common
auxiliary vocabulary and sentences \(\widehat\varphi_n(d)\).
Add a fresh unary predicate \(I\) and a fresh binary relation \(<\).
Let \(\lambda\) be the first-order sentence saying that \(<\) relates
only elements of \(I\) and is a strict linear order on \(I\).
It does not require \(I\) to be nonempty. For \(n<\omega\), let
\[
 H_n(d)\ :=\
 \exists y_0\cdots\exists y_n
 \left(\bigwedge_{j\leq n}(I(y_j)\land y_j<d)
       \land\bigwedge_{j<\ell\leq n} y_j\ne y_\ell\right).
\]
It states that \(d\) has at least \(n+1\) predecessors in \(I\).
Take the theory
\[
 \Psi=\{\lambda\}\ \cup\
 \left\{\forall d\bigl((I(d)\land H_n(d))\to
                    \widehat\varphi_n(d)\bigr):n<\omega\right\}.
\]
Boolean closure and particularization put every sentence in
\(\mathscr L\). This theory is countable even when \(\tau\) and the
auxiliary vocabulary are uncountable. Its cardinality is the number
of its sentences, not the number of its symbols.

In a model \(\mathfrak B\models\Psi\), an element \(b\in I\) with
infinitely many predecessors would satisfy every \(H_n\). The single
induced structure \(\mathfrak B_b\) would then satisfy every
\(\varphi_n\), by the fibre theorem. This contradicts the choice of
\(T\). Thus each element of \(I\) has finitely many predecessors.
To see that \((I,<)\) is well-ordered, take a nonempty subset \(X\)
and \(b\in X\). The elements of \(X\) at or below \(b\) form a
nonempty finite linear order. Its least member is least in \(X\),
since every other member of \(X\) lies above \(b\).
Moreover, sending an element to its number of predecessors identifies
\((I,<)\) with an initial segment of \(\omega\): this map is strictly
increasing, and an element with \(m\) predecessors has exactly one
predecessor of each rank below \(m\). Thus the order type is finite
or \(\omega\). The field of \(<\) is \(I\) if \(I\) has at least
two elements and is empty otherwise. It is therefore always well-ordered.

For each \(n<\omega\), choose a model \(\mathfrak M_n\) of
\(\{\varphi_j:j<n\}\). For \(n=0\), choose any nonempty
\(\tau\)-structure. The fibre theorem realizes these models on
disjoint fibres over indices \(i_n\). Interpret \(I\) as these
indices, and put \(i_j<i_n\) exactly when \(j<n\), with no other
ordered pairs. The index \(i_n\) has exactly \(n\) predecessors.
Hence \(H_j(i_n)\) holds exactly when \(j<n\), in which case
\(\mathfrak M_n\models\varphi_j\). The fibre equivalence proves
every conditional in \(\Psi\). This is a model of \(\Psi\) whose
ordered field has type \(\omega\); consequently \(\Psi\) pins
\(\omega\).

Since the well-ordering number is a supremum of the pinned ordinals,
we also construct a pinned ordinal strictly above \(\omega\).
Add a fresh constant \(e\), the sentence \(\neg I(e)\), a fresh
binary relation \(\prec\), and its first-order definition
\[
 \forall u\,\forall v\,
 \bigl(u\prec v\leftrightarrow
       (u<v\lor(I(u)\land v=e))\bigr).
\]
In every model this puts \(e\) after all of \(I\). If \(I\) is empty,
its field is empty; otherwise its field is \(I\cup\{e\}\), ordered
as the well-order on \(I\) with a final element. All models therefore
well-order the field. In the constructed \(\omega\)-model, interpret
\(e\) by any element of \(M_0\), which is nonempty and disjoint from
\(I\). This gives order type \(\omega+1\).
The enlarged theory is still countable, so
\(w_{\aleph_0}(\mathscr L)\geq\omega+1>\omega\).
We have proved the contrapositive of (3) implies (1), and also both
additional pinning assertions. Together with the first two implications,
this proves the equivalence.
\end{proof}
\end{annexitem}
